% ============================================================
% issuer.tex — Aussteller: wer abrechnet, wo er sitzt, wie man ihn bezahlt.
%
% Hier stehen bewusst NUR erfundene Musterdaten, identisch mit
% stammdaten.example/ich.yaml. Sie dienen dem Beispiel-PDF (vorlage/beispiel.pdf).
% Für echte Rechnungen erzeugt render.py diese Datei zur Laufzeit im Temp-Verzeichnis
% aus rechnung.resolved.json (Vorlage: issuer.tex.j2); die echten Daten liegen im
% Gewerbe-Ordner und niemals in diesem Repo.
% ============================================================

% ---------- Briefkopf (Seite 1, unter der Rahmenlinie) ----------
\newcommand{\InvIssuerName}{Mustermann Beratung}           % fette Zeile
\newcommand{\InvIssuerPerson}{Max Mustermann}              % natürliche Person: Gruß, Kontoinhaber
\newcommand{\InvIssuerAddressLine}{Musterstraße 1 \textperiodcentered\ 12345 Musterstadt}
\newcommand{\InvIssuerCountryEn}{Germany}                  % an die Adresszeile in englischen Rechnungen
\newcommand{\InvIssuerCountryDe}{}                         % … und in deutschen ({} = weglassen, Inland)
\newcommand{\InvIssuerEmail}{max@example.com}
\newcommand{\InvIssuerPhone}{+49 123 4567890}
\newcommand{\InvIssuerWeb}{example.com}

% ---------- Rücksendeangabe (klein über der Empfängeranschrift) ----------
% Nur Name, Straße, Ort. Keine Kontaktdaten – sie stünden sonst im Fenster des Umschlags.
\newcommand{\InvIssuerSenderLine}{Max Mustermann \textperiodcentered\ Musterstraße 1 \textperiodcentered\ 12345 Musterstadt}

% ---------- Steuerliche Kennungen (Metadatenblock und Fußzeile) ----------
\newcommand{\InvIssuerTaxNo}{143/123/12345}                % Steuernummer
\newcommand{\InvIssuerVatId}{DE123456789}                  % USt-IdNr.

% ---------- Bank (Zahlungssatz und Fußzeile) ----------
\newcommand{\InvIssuerIban}{DE89 3704 0044 0532 0130 00}
\newcommand{\InvIssuerBic}{COBADEFFXXX}
\newcommand{\InvIssuerBank}{Musterbank}
